\section{Conclusion and Future Work}
\label{sec:conclusion}

This paper studied trace-complexity bias in a common microservice
anomaly scoring pattern: Service Trace Vector representations scored by
an unnormalized unsupervised detector. On real, non-injected TrainTicket
traces, baseline STV~+~Isolation Forest anomaly scores correlate strongly
with trace size, with correlations of approximately
\(r \approx 0.85\)--\(0.98\) for span and service-count based complexity
measures. This finding shows that raw STV-based scores can reflect trace
structural complexity rather than only operational abnormality.

To address this issue, we proposed WR-STV, a lightweight STV-based
scoring method that applies per-feature residual standardization,
residual clipping, reliability weighting, and top-\(k\) aggregation.
Across the evaluated datasets, WR-STV substantially reduces the measured
correlation between anomaly score and trace complexity
(\(r \approx 0.30\)) while improving controlled STV-visible latency
anomaly ranking over baseline STV~+~Isolation Forest
(within-dataset ROC-AUC \(0.605 \rightarrow 0.814\)). Cross-dataset
experiments also show promising transfer behavior
(\(0.622 \rightarrow 0.803\)), although these results should be
interpreted cautiously because they are based on only six directed
train-test pairs.

The ablation results indicate that most of the improvement comes from
feature-wise residual normalization, while the proposed reliability
weighting contributes a smaller, metric-dependent refinement. WR-STV
also provides interpretable service-path-level evidence, with top-3
contributing service paths recovering more than \(90\%\) of known
injected paths in the evaluated setting. Runtime analysis further shows
that WR-STV achieves roughly a \(16\times\) lower per-trace scoring time
than Isolation Forest-based baselines. Overall, the results support
WR-STV as a lightweight and interpretable approach for reducing
trace-complexity bias in controlled STV-visible latency anomaly scoring,
rather than as a general solution to all microservice anomaly types.

\subsection*{Future Work}

Future work will extend this study in four directions. First, evaluation
should include anomalies beyond STV-visible latency perturbations,
including dropped calls, novel call paths, missing spans, and other
structural trace anomalies. Second, WR-STV should be evaluated on
additional services, benchmark systems, and production-like trace
datasets to test whether the observed complexity-bias pattern
generalizes beyond \texttt{ts-order-service}. Third, the failure analysis
suggests that very short traces are a predictable low-evidence case;
future systems could attach confidence estimates to anomaly scores or
abstain from scoring traces below a minimum evidence threshold. Finally,
WR-STV should be evaluated against real fault-injection or incident-labeled
datasets so that its behavior can be assessed on anomalies not generated
by the same evaluation procedure used in this paper.